\documentclass[11pt]{article}
\usepackage[margin=1in]{geometry}
\usepackage{booktabs}
\usepackage{array}
\usepackage{tabularx}
\usepackage{enumitem}
\usepackage[T1]{fontenc}
\usepackage[utf8]{inputenc}
\usepackage{lmodern}
\setlength{\parskip}{0.6em}
\setlength{\parindent}{0pt}
\setlist[itemize]{topsep=3pt,itemsep=3pt,leftmargin=1.4em}
\setlist[enumerate]{topsep=3pt,itemsep=3pt,leftmargin=1.6em}
\title{Redwood Regional Bank CECL Review Memo}
\date{March 13, 2026}
\begin{document}
\maketitle
\section*{Executive Summary}
This memo documents an end-to-end CECL review of Redwood Regional Bank's Q1 2026 CECL Allowance Review. The supplied reserve engine was executable, the packaged baseline reserve was reproduced exactly, and scenario reruns plus targeted sensitivities were completed using the supplied portfolio, scenario, and overlay artifacts. Portfolio-level reserve direction was broadly reasonable with total ACL rising from \$4,876,841 under Baseline to \$7,132,022 under Adverse and \$8,287,703 under Severe. The review nonetheless identified material exceptions requiring remediation before an unqualified approval posture would be appropriate.
\section*{Scope and Coverage}
This review assessed the supplied CECL package for internal coherence, quantitative reasonableness, and documentation alignment. The work covered baseline reproduction, scenario reruns, sensitivity testing, reserve driver analysis, and cross-checking of methodology, scenario, and overlay documentation against implemented behavior.
\begin{itemize}
\item Supported work: baseline reproduction, scenario reruns, sensitivity testing, driver analysis, and documentation cross-checking.
\item Blocked work: none within the scope of the supplied package.
\item Coverage note: conclusions are grounded in the supplied engine, portfolio, scenario files, overlays, and supporting documentation.
\end{itemize}
\section*{Artifacts Reviewed}
\begin{itemize}
\item model/cecl\_engine.py (code)
\item data/loan\_level\_snapshot.csv (dataset)
\item data/data\_dictionary.csv (document)
\item data/overlay\_schedule.csv (config)
\item scenarios/baseline.csv (scenario)
\item scenarios/adverse.csv (scenario)
\item scenarios/severe.csv (scenario)
\item docs/methodology.md (document)
\item docs/model\_overview.md (document)
\item docs/scenario\_assumptions.md (document)
\item docs/overlay\_memo.md (document)
\item outputs/prior\_baseline\_reserve.csv (metrics)
\end{itemize}
\section*{Package Discovery and Case Understanding}
The discovered package is best understood as a model-driven CECL review case for Redwood Regional Bank's Q1 2026 CECL Allowance Review. It includes executable reserve logic, loan-level portfolio data, scenario definitions, prior baseline outputs, and enough supporting documentation to support both quantitative reruns and documentation challenge.
\begin{itemize}
\item Central assumption or theme: Reasonable-and-supportable horizon is documented as 4 quarters, with 4 quarters of reversion.
\item Central assumption or theme: Scenario severity should increase reserve directionally from Baseline to Adverse to Severe at portfolio level and generally by segment.
\item Central assumption or theme: Qualitative overlays should remain within documented governance guardrails and should not produce unexplained segment relief.
\item Reviewable scope: Baseline reproduction from supplied engine and data.
\item Reviewable scope: Scenario reruns and segment-level reasonableness testing.
\item Reviewable scope: Sensitivity testing for horizon, reversion, macro severity, and overlays.
\item Reviewable scope: Methodology, scenario, and overlay documentation alignment review.
\item Key risk to challenge: Documentation may not match the implemented forecast and reversion treatment.
\item Key risk to challenge: Overlay governance may not reflect actual overlay usage by scenario and segment.
\item Key risk to challenge: A severe scenario can appear directionally harsher at portfolio level while still producing a segment anomaly.
\item Constraint: No independent historical performance dataset or external benchmark was supplied; conclusions are bounded to the package and synthetic portfolio.
\end{itemize}
\section*{CECL Process Understanding}
The package describes a loan-level CECL process spanning 4 segments: Commercial Real Estate, Commercial and Industrial, Residential Mortgage, Consumer Unsecured. The quantitative framework applies scenario-dependent macro paths to lifetime PD and LGD components and then layers qualitative overlays by segment and scenario.
\begin{center}
\small
\begin{tabularx}{\textwidth}{>{\raggedright\arraybackslash}p{0.171\textwidth}>{\raggedright\arraybackslash}p{0.235\textwidth}>{\raggedright\arraybackslash}p{0.256\textwidth}>{\raggedright\arraybackslash}p{0.299\textwidth}}
\toprule
Assumption Area & Documented Position & Observed in Package & Review Implication \\
\midrule
Forecast horizon & 4 quarters documented & 6 quarters implemented & Methodology and model overview are not aligned to engine behavior. \\
Reversion horizon & 4 quarters documented & 2 quarters implemented & Policy description and actual reserve mechanics differ. \\
Overlay cap & 5.0 bps documented & 9.0 bps scheduled & Qualitative adjustment magnitude exceeds documented guardrail. \\
Scenario reasonableness & Severe should be directionally harsher unless clearly justified & Residential Mortgage reserve falls by \$109,988 from adverse to severe & Segment-level severe path requires additional governance challenge. \\
\bottomrule
\end{tabularx}
\normalsize
\end{center}
\section*{Review Strategy and Planning Logic}
The review strategy was formed from the discovered evidence. Codex considered 5 central review questions, selected 10 executable procedures, and identified 0 blocked or non-applicable procedures.
Because the package contained runnable code, scenario files, prior outputs, and core methodology documents, the review strategy prioritized execution-based challenge first and then used documentation review to test whether the written methodology and governance artifacts actually described the observed reserve behavior.
\begin{enumerate}
\item Does the package contain enough evidence to support a full execution-based CECL review rather than a documentation-only assessment?
\item Can the supplied baseline reserve be reproduced from the packaged reserve engine, portfolio data, and baseline scenario?
\item Do reserve outputs move directionally and materially across Baseline, Adverse, and Severe scenarios both overall and by major segment?
\item How sensitive are reserves to forecast horizon, reversion speed, macro severity, and overlay magnitude assumptions?
\item Do the methodology, scenario, and overlay documents faithfully describe the behavior observed in the implemented CECL process, especially for Residential Mortgage?
\end{enumerate}
\begin{center}
\small
\begin{tabularx}{\textwidth}{>{\raggedright\arraybackslash}p{0.107\textwidth}>{\raggedright\arraybackslash}p{0.235\textwidth}>{\raggedright\arraybackslash}p{0.149\textwidth}>{\raggedright\arraybackslash}p{0.469\textwidth}}
\toprule
ID & Procedure & Status & Why Selected \\
\midrule
FR-01 & Package inventory and case classification & Executed & Runnable code, scenario tables, loan-level data, prior outputs, and methodology documents were all present. \\
FR-02 & Baseline reproduction & Executed & A packaged prior baseline reserve was supplied alongside the runnable engine and source portfolio. \\
FR-03 & Portfolio-level scenario reruns & Executed & The package included three explicit scenario files and a model designed to consume them. \\
FR-04 & Segment-level reasonableness review & Executed & CECL governance depends on segment-level reserve understanding, not just portfolio totals. \\
FR-05 & Forecast horizon sensitivity & Executed & The methodology makes forecast horizon a named governance assumption. \\
FR-06 & Reversion speed sensitivity & Executed & Reversion mechanics are central to CECL governance and are explicitly documented in the package. \\
FR-07 & Macro severity sensitivity & Executed & Scenario severity and directional reasonableness are core to CECL challenge work. \\
FR-08 & Overlay magnitude sensitivity & Executed & The package included a scenario-by-segment overlay schedule and an overlay governance memo. \\
FR-09 & Reserve driver bridge & Executed & A driver bridge makes the review more interpretable and supports reasonableness challenge. \\
FR-10 & Methodology and documentation alignment review & Executed & The package contained enough documentation to support direct cross-checking against the model and results. \\
\bottomrule
\end{tabularx}
\normalsize
\end{center}
\section*{Procedures Performed}
\subsection*{FR-01 - Package inventory and case classification}
Determine whether the upload supports a model-driven CECL review and identify the core evidence types available.
\begin{itemize}
\item Status: executed
\item Why selected: Runnable code, scenario tables, loan-level data, prior outputs, and methodology documents were all present.
\item Evidence relied upon: model/cecl\_engine.py, data/loan\_level\_snapshot.csv, scenarios/baseline.csv, outputs/prior\_baseline\_reserve.csv, docs/methodology.md
\item Outputs produced: discovery\_summary.md, evidence\_ledger.json
\item Assessment criteria: Review should only proceed as execution-based if engine, data, scenarios, and prior outputs are all available.
\item Key result: Package classified as model-driven CECL review with no material execution blockers.
\end{itemize}
\subsection*{FR-02 - Baseline reproduction}
Confirm that the packaged baseline reserve can be reproduced from the supplied engine and portfolio inputs.
\begin{itemize}
\item Status: executed
\item Why selected: A packaged prior baseline reserve was supplied alongside the runnable engine and source portfolio.
\item Evidence relied upon: model/cecl\_engine.py, data/loan\_level\_snapshot.csv, outputs/prior\_baseline\_reserve.csv
\item Outputs produced: baseline\_reproduction.json
\item Assessment criteria: Total reserve and segment reserve should reconcile within tolerance; large deltas would undermine confidence in package integrity.
\item Key result: Baseline rerun matched packaged reserve exactly at 4876841.47 with 0.000\% delta.
\end{itemize}
\subsection*{FR-03 - Portfolio-level scenario reruns}
Run Baseline, Adverse, and Severe scenarios and assess directional reserve ordering.
\begin{itemize}
\item Status: executed
\item Why selected: The package included three explicit scenario files and a model designed to consume them.
\item Evidence relied upon: scenarios/baseline.csv, scenarios/adverse.csv, scenarios/severe.csv, model/cecl\_engine.py
\item Outputs produced: scenario\_results.csv
\item Assessment criteria: Portfolio reserve should generally increase with scenario severity absent a documented exception.
\item Key result: Portfolio reserve increased from 4876841.47 to 7132021.97 to 8287703.47.
\end{itemize}
\subsection*{FR-04 - Segment-level reasonableness review}
Test whether scenario ordering remains directionally sensible at the segment level.
\begin{itemize}
\item Status: executed
\item Why selected: CECL governance depends on segment-level reserve understanding, not just portfolio totals.
\item Evidence relied upon: outputs/support/segment\_reserve\_comparison.csv, docs/scenario\_assumptions.md, docs/overlay\_memo.md
\item Outputs produced: segment\_reserve\_comparison.csv
\item Assessment criteria: Segment movements should generally align to scenario severity or be explicitly supported by documented assumptions.
\item Key result: Residential Mortgage severe reserve was 217710.91, below adverse reserve 327699.21, creating a documented challenge item.
\end{itemize}
\subsection*{FR-05 - Forecast horizon sensitivity}
Assess reserve sensitivity to changing the reasonable-and-supportable forecast horizon.
\begin{itemize}
\item Status: executed
\item Why selected: The methodology makes forecast horizon a named governance assumption.
\item Evidence relied upon: docs/methodology.md, outputs/support/sensitivity\_results.csv
\item Outputs produced: sensitivity\_results.csv
\item Assessment criteria: Reserve movement should be understandable and documented; mismatch to stated policy must be surfaced.
\item Key result: Moving from 6 to 4 forecast quarters changed baseline reserve only modestly, but it directly exposed a documented-versus-implemented policy mismatch.
\end{itemize}
\subsection*{FR-06 - Reversion speed sensitivity}
Assess how severe-scenario reserve responds to longer reversion periods.
\begin{itemize}
\item Status: executed
\item Why selected: Reversion mechanics are central to CECL governance and are explicitly documented in the package.
\item Evidence relied upon: docs/methodology.md, outputs/support/sensitivity\_results.csv
\item Outputs produced: sensitivity\_results.csv
\item Assessment criteria: Reserve should move directionally as reversion assumptions are extended or shortened.
\item Key result: Extending severe-scenario reversion from 2 to 6 quarters increased reserve materially, confirming reversion is a meaningful reserve lever.
\end{itemize}
\subsection*{FR-07 - Macro severity sensitivity}
Test whether reserve responds proportionally to macro severity scaling.
\begin{itemize}
\item Status: executed
\item Why selected: Scenario severity and directional reasonableness are core to CECL challenge work.
\item Evidence relied upon: scenarios/severe.csv, outputs/support/sensitivity\_results.csv
\item Outputs produced: sensitivity\_results.csv
\item Assessment criteria: Scaling macro severity upward should generally increase reserve; weaker response would require explanation.
\item Key result: Scaling severe stress from 1.00x to 1.25x materially increased reserve, supporting macro sensitivity monotonicity at portfolio level.
\end{itemize}
\subsection*{FR-08 - Overlay magnitude sensitivity}
Assess the reserve effect of removing or amplifying qualitative overlays.
\begin{itemize}
\item Status: executed
\item Why selected: The package included a scenario-by-segment overlay schedule and an overlay governance memo.
\item Evidence relied upon: data/overlay\_schedule.csv, docs/overlay\_memo.md, outputs/support/sensitivity\_results.csv
\item Outputs produced: sensitivity\_results.csv
\item Assessment criteria: Overlay effect should be quantitatively explainable and consistent with governance statements.
\item Key result: Removing severe overlays reduced reserve materially, confirming overlay posture is consequential and governance-relevant.
\end{itemize}
\subsection*{FR-09 - Reserve driver bridge}
Identify which macro and overlay drivers explain the majority of severe-versus-baseline reserve movement.
\begin{itemize}
\item Status: executed
\item Why selected: A driver bridge makes the review more interpretable and supports reasonableness challenge.
\item Evidence relied upon: outputs/support/driver\_bridge.csv, scenarios/baseline.csv, scenarios/severe.csv
\item Outputs produced: driver\_bridge.csv
\item Assessment criteria: Major reserve movement should be attributable to understandable drivers rather than unexplained residual behavior.
\item Key result: CRE price growth was the largest modeled standalone driver of severe-versus-baseline reserve change.
\end{itemize}
\subsection*{FR-10 - Methodology and documentation alignment review}
Test whether written methodology, model overview, and overlay documentation match implemented behavior and observed outputs.
\begin{itemize}
\item Status: executed
\item Why selected: The package contained enough documentation to support direct cross-checking against the model and results.
\item Evidence relied upon: docs/methodology.md, docs/model\_overview.md, docs/scenario\_assumptions.md, docs/overlay\_memo.md, model/cecl\_engine.py
\item Outputs produced: documentation\_crosscheck.md, findings\_register.json
\item Assessment criteria: Documented horizon, reversion, scenario logic, and overlay posture should match the implemented process or be clearly qualified.
\item Key result: Cross-checking identified a forecast/reversion mismatch, overlay cap mismatch, and a severe-scenario segment anomaly requiring remediation.
\end{itemize}
\section*{Agentic Execution Workflow}
\subsection*{Discovery}
The review started by inventorying the package and determining whether a quantitative review was actually possible from the supplied materials.
\begin{itemize}
\item Review question addressed: What is in the package and does it support execution-based CECL review?
\item Decision or rationale: Treat the case as model-driven because reserve code, loan-level data, scenarios, prior outputs, and core methodology documents were all present.
\item Inputs reviewed: input\_package/*
\item Outputs produced: discovery\_summary.md, evidence\_ledger.json, case\_understanding.md
\end{itemize}
\subsection*{Case understanding and planning}
Codex formed a review strategy around the strongest available evidence and explicitly selected a test set before running the work.
\begin{itemize}
\item Review question addressed: Which review questions matter most for this package and which procedures can answer them credibly?
\item Decision or rationale: Prioritize baseline reproduction, scenario reruns, sensitivities, and documentation alignment because the package supports all four lanes.
\item Inputs reviewed: discovery\_summary.md, docs/methodology.md, docs/scenario\_assumptions.md, docs/overlay\_memo.md
\item Outputs produced: review\_plan.md, review\_strategy.md, executed\_test\_matrix.md
\end{itemize}
\subsection*{Quantitative execution}
The review ran the CECL engine against baseline and stressed scenarios, then executed targeted sensitivities to challenge the major reserve assumptions.
\begin{itemize}
\item Review question addressed: Can the reserve process be reproduced, and do outputs move sensibly overall and by segment?
\item Decision or rationale: Run baseline, adverse, and severe scenarios first, then test horizon, reversion, macro severity, and overlays as the most consequential levers.
\item Inputs reviewed: model/cecl\_engine.py, data/loan\_level\_snapshot.csv, scenarios/*.csv, data/overlay\_schedule.csv
\item Outputs produced: baseline\_reproduction.json, scenario\_results.csv, segment\_reserve\_comparison.csv, sensitivity\_results.csv, driver\_bridge.csv
\end{itemize}
\subsection*{Documentation challenge}
After quantitative outputs were produced, the review cross-checked the documents against implemented behavior and reserve results.
\begin{itemize}
\item Review question addressed: Do the written methodology, scenario, and overlay materials describe the same process that the results imply?
\item Decision or rationale: Escalate mismatches where documented governance assumptions differ from implementation or where narrative support is insufficient for observed output behavior.
\item Inputs reviewed: docs/*.md, model/cecl\_engine.py, segment\_reserve\_comparison.csv, sensitivity\_results.csv
\item Outputs produced: documentation\_crosscheck.md, findings\_register.json, evidence\_map.md
\end{itemize}
\subsection*{Synthesis}
The final stage combined discovery, planning, quantitative results, and documentation challenge into an internal review report and support pack.
\begin{itemize}
\item Review question addressed: What opinion can be supported, what findings should be issued, and what evidence supports that conclusion?
\item Decision or rationale: Issue a substantive review with explicit remediation for documentation mismatch, overlay governance, and the residential mortgage severe-scenario anomaly.
\item Inputs reviewed: findings\_register.json, executed\_test\_matrix.md, coverage\_statement.md, agentic\_review\_log.md
\item Outputs produced: cecl\_review\_memo.tex, cecl\_review\_memo.pdf
\end{itemize}
\section*{Quantitative Results}
\subsection*{Baseline Reproduction}
Baseline reproduction landed at \$4,876,841 against a packaged baseline of \$4,876,841, with absolute delta \$0 and maximum loan-level delta \$0. Relative to Baseline, Adverse increased reserve by \$2,255,180 and Severe increased reserve by another \$1,155,682, for a total Severe-versus-Baseline increase of \$3,410,862.
\begin{center}
\begin{tabularx}{\textwidth}{>{\raggedright\arraybackslash}p{0.538\textwidth}>{\raggedright\arraybackslash}p{0.422\textwidth}}
\toprule
Metric & Result \\
\midrule
Packaged baseline reserve & \$4,876,841 \\
Rerun baseline reserve & \$4,876,841 \\
Absolute delta & \$0 \\
Percent delta & 0.000\% \\
Maximum loan-level absolute delta & \$0 \\
\bottomrule
\end{tabularx}
\end{center}
\subsection*{Scenario Results}
Portfolio-level reserve ordering was monotonic. Adverse increased reserve by \$2,255,180 versus Baseline, and Severe increased reserve by \$1,155,682 versus Adverse.
\begin{center}
\small
\begin{tabularx}{\textwidth}{>{\raggedright\arraybackslash}p{0.216\textwidth}>{\raggedright\arraybackslash}p{0.264\textwidth}>{\raggedright\arraybackslash}p{0.216\textwidth}>{\raggedright\arraybackslash}p{0.264\textwidth}}
\toprule
Scenario & Reserve & Reserve Rate & Change vs Baseline \\
\midrule
Baseline & \$4,876,841 & 0.60\% & \$0 \\
Adverse & \$7,132,022 & 0.88\% & \$2,255,180 \\
Severe & \$8,287,703 & 1.03\% & \$3,410,862 \\
\bottomrule
\end{tabularx}
\normalsize
\end{center}
\subsection*{Segment Results}
Segment-level behavior was directionally reasonable for most segments. The notable exception was Residential Mortgage, which declined by \$109,988 from Adverse to Severe.
\begin{center}
\small
\begin{tabularx}{\textwidth}{>{\raggedright\arraybackslash}p{0.196\textwidth}>{\raggedright\arraybackslash}p{0.153\textwidth}>{\raggedright\arraybackslash}p{0.153\textwidth}>{\raggedright\arraybackslash}p{0.153\textwidth}>{\raggedright\arraybackslash}p{0.153\textwidth}>{\raggedright\arraybackslash}p{0.153\textwidth}}
\toprule
Segment & Baseline & Adverse & Severe & Adv - Base & Sev - Adv \\
\midrule
Commercial And Industrial & \$1,977,246 & \$2,712,809 & \$3,148,324 & \$735,563 & \$435,515 \\
Commercial Real Estate & \$2,545,520 & \$3,972,723 & \$4,787,770 & \$1,427,204 & \$815,047 \\
Consumer Unsecured & \$100,356 & \$118,791 & \$133,899 & \$18,435 & \$15,108 \\
Residential Mortgage & \$253,720 & \$327,699 & \$217,711 & \$73,979 & \$-109,988 \\
\bottomrule
\end{tabularx}
\normalsize
\end{center}
\subsection*{Sensitivity Testing}
Sensitivity testing showed limited change from moving the documented forecast horizon to 4 quarters (change of \$1,145), but reserve meaningfully increased as reversion lengthened from 2 to 6 quarters under Severe (change of \$185,201). Scaling severe macro stress from 1.00x to 1.25x added \$1,515,402, while removing overlays reduced Severe reserve by \$436,787.
\begin{center}
\small
\begin{tabularx}{\textwidth}{>{\raggedright\arraybackslash}p{0.216\textwidth}>{\raggedright\arraybackslash}p{0.216\textwidth}>{\raggedright\arraybackslash}p{0.264\textwidth}>{\raggedright\arraybackslash}p{0.264\textwidth}}
\toprule
Forecast Horizon & Scenario & Reserve & Change vs 6Q Base \\
\midrule
4 quarters & Baseline & \$4,877,987 & \$1,145 \\
6 quarters & Baseline & \$4,876,841 & \$0 \\
8 quarters & Baseline & \$4,876,841 & \$0 \\
\bottomrule
\end{tabularx}
\normalsize
\end{center}
\begin{center}
\small
\begin{tabularx}{\textwidth}{>{\raggedright\arraybackslash}p{0.216\textwidth}>{\raggedright\arraybackslash}p{0.216\textwidth}>{\raggedright\arraybackslash}p{0.264\textwidth}>{\raggedright\arraybackslash}p{0.264\textwidth}}
\toprule
Reversion Speed & Scenario & Reserve & Change vs 2Q Base \\
\midrule
2 quarters & Severe & \$8,287,703 & \$0 \\
4 quarters & Severe & \$8,384,331 & \$96,627 \\
6 quarters & Severe & \$8,472,904 & \$185,201 \\
\bottomrule
\end{tabularx}
\normalsize
\end{center}
\begin{center}
\small
\begin{tabularx}{\textwidth}{>{\raggedright\arraybackslash}p{0.216\textwidth}>{\raggedright\arraybackslash}p{0.216\textwidth}>{\raggedright\arraybackslash}p{0.264\textwidth}>{\raggedright\arraybackslash}p{0.264\textwidth}}
\toprule
Severity Scale & Scenario & Reserve & Change vs 1.00x \\
\midrule
0.75x & Severe & \$7,126,890 & \$-1,160,813 \\
1.0x & Severe & \$8,287,703 & \$0 \\
1.25x & Severe & \$9,803,106 & \$1,515,402 \\
\bottomrule
\end{tabularx}
\normalsize
\end{center}
\begin{center}
\small
\begin{tabularx}{\textwidth}{>{\raggedright\arraybackslash}p{0.216\textwidth}>{\raggedright\arraybackslash}p{0.216\textwidth}>{\raggedright\arraybackslash}p{0.264\textwidth}>{\raggedright\arraybackslash}p{0.264\textwidth}}
\toprule
Overlay Multiplier & Scenario & Reserve & Change vs 1.00x \\
\midrule
0.0x & Severe & \$7,850,917 & \$-436,787 \\
1.0x & Severe & \$8,287,703 & \$0 \\
1.5x & Severe & \$8,506,097 & \$218,393 \\
\bottomrule
\end{tabularx}
\normalsize
\end{center}
\subsection*{Driver Analysis}
The simple reserve bridge indicates that Cre Price Growth was the largest modeled contributor to the Severe-versus-Baseline reserve increase, adding \$1,108,566 on a standalone basis.
\begin{center}
\begin{tabularx}{\textwidth}{>{\raggedright\arraybackslash}p{0.538\textwidth}>{\raggedright\arraybackslash}p{0.422\textwidth}}
\toprule
Driver & Reserve Delta vs Baseline \\
\midrule
Cre Price Growth & \$1,108,566 \\
Gdp Growth & \$455,249 \\
Overlay & \$436,787 \\
Unemployment Rate & \$336,020 \\
House Price Growth & \$169,118 \\
Prime Rate & \$66,886 \\
\bottomrule
\end{tabularx}
\end{center}
\section*{Documentation and Assumption Alignment}
Documentation cross-checking found that the written methodology still describes a 4-quarter reasonable-and-supportable period with 4-quarter reversion, while the engine implements 6 forecast quarters and 2 reversion quarters. The overlay memo references a 5.0 bps guardrail, but the supplied overlay schedule applies up to 9.0 bps. At segment level, Residential Mortgage shows a severe-scenario reserve \$217,711, below its adverse reserve \$327,699, driven by the scenario and overlay posture described in the package.
\begin{itemize}
\item Documented forecast/reversion: 4Q forecast and 4Q reversion.
\item Implemented forecast/reversion: 6Q forecast and 2Q reversion.
\item Documented overlay cap: 5.0 bps; scheduled overlay cap: 9.0 bps.
\item Anomaly observation: Residential Mortgage severe reserve \$217,711 versus adverse reserve \$327,699.
\end{itemize}
\section*{Detailed Findings and Remediation}
\subsection*{Documented CECL horizon and reversion do not match the implemented reserve engine}
\textbf{Severity:} HIGH\\
\textbf{Observation:} Documentation describes a 4-quarter reasonable-and-supportable period with 4-quarter reversion, while the reserve engine is configured for 6 forecast quarters and 2 reversion quarters.\\
\textbf{Evidence:} docs/methodology.md, docs/model\_overview.md, model/cecl\_engine.py\\
\textbf{Recommended response:} Either update the methodology and model overview to the implemented 6-quarter forecast and 2-quarter reversion, or change the engine to match documented policy before the next governance submission.
\subsection*{Qualitative overlay effect is materially larger than documentation suggests}
\textbf{Severity:} MEDIUM\\
\textbf{Observation:} The documentation frames qualitative overlays as capped at 5.0 bps, but the supplied overlay table uses up to 9.0 bps.\\
\textbf{Evidence:} docs/overlay\_memo.md, data/overlay\_schedule.csv, outputs/support/sensitivity\_results.csv\\
\textbf{Recommended response:} Reconcile the overlay memo and governance materials to the actual schedule, including the 9.0 bps cap and the rationale for segment-level relief under Severe.
\subsection*{One segment behaves directionally oddly under the severe scenario}
\textbf{Severity:} HIGH\\
\textbf{Observation:} Segment residential\_mortgage produces a severe-scenario reserve of 217710.91, below the adverse reserve of 327699.21. The severe path allows house prices to recover faster than adverse after 2026Q4 while the overlay schedule applies negative mortgage relief, creating a directionally odd segment result.\\
\textbf{Evidence:} outputs/support/segment\_reserve\_comparison.csv, scenarios/severe.csv, docs/scenario\_assumptions.md, docs/overlay\_memo.md\\
\textbf{Recommended response:} Challenge the Residential Mortgage severe scenario result, document whether the house-price recovery and negative overlay are intended, and update committee materials accordingly.
\subsection*{Baseline reserve was reproducible within tolerance}
\textbf{Severity:} LOW\\
\textbf{Observation:} Rerunning the supplied baseline logic reproduced the packaged baseline reserve within 0.000\%.\\
\textbf{Evidence:} outputs/prior\_baseline\_reserve.csv, outputs/support/baseline\_reproduction.json\\
\textbf{Recommended response:} Retain the reproduction evidence as support for package completeness and execution readiness.
\section*{Overall Assessment}
The supplied package supported a substantive CECL review and allowed execution of the core quantitative and documentation work expected for this scope. The baseline output was reproducible and overall reserve direction was reasonable, but documentation and governance artifacts are not yet aligned to implemented behavior, and one segment-level severe result requires explicit management challenge before the package would support an unqualified review conclusion.
\end{document}